% \section{Introduction}

% Urban air pollution is shaped by a combination of local emissions,
% meteorological transport, dispersion, regional background pollution, and
% sensor noise. A common goal in air-quality monitoring is source
% apportionment: estimating how much different source categories, such as
% traffic, industry, brick kilns, road dust, or population-related activity,
% contribute to observed pollution levels. This problem is especially
% important for policy, since mitigation decisions depend not only on
% predicting pollution concentrations, but on attributing those
% concentrations to actionable source groups.

% A simplified description of city-scale pollutant transport is
% \begin{align}
% \partial_t u(\mathbf x,t)
% &=
% \mathcal F_\eta
% \big(
% u(\mathbf x,t),
% \nabla u(\mathbf x,t),
% \Delta u(\mathbf x,t);
% \mathbf w(t)
% \big)
% \nonumber \\
% &
% \qquad
% +
% s(\mathbf x,t)
% +
% b(\mathbf x,t),
% \nonumber \\
% &\quad
% \mathbf x\in\Omega\subset\mathbb R^d,\quad t\in[0,T],
% \end{align}
% where \(u(\mathbf x,t)\) denotes pollutant concentration,
% \(\mathbf w(t)\) denotes wind and meteorological conditions,
% \(s(\mathbf x,t)\) denotes local emissions, \(b(\mathbf x,t)\) denotes
% background or regional pollution, and \(\eta\) represents transport and
% dispersion parameters. In practice, however, the concentration field is
% not fully observed. Measurements are available only at a small number of
% fixed monitoring stations, producing sparse and noisy sensor readings.

% This sparse sensing regime makes source apportionment fundamentally
% different from standard prediction. A model may accurately reconstruct
% sensor trajectories while still failing to uniquely identify the
% underlying source contributions. The difficulty is not only that
% transport models are approximate or that emissions inventories are
% imperfect. Even with a fixed transport model and known candidate source
% inventories, different source groups may induce nearly indistinguishable
% sensor-time signatures under the available wind and sensor geometry.

% Rather than attempting to recover an arbitrary unknown source field, we
% study inventory-based source apportionment. Let
% \[
% S
% =
% \begin{bmatrix}
% \mathbf s_1 &
% \mathbf s_2 &
% \cdots &
% \mathbf s_K
% \end{bmatrix}
% \in
% \mathbb R^{q\times K}
% \]
% denote a collection of known or proxy source maps, where each column
% \(\mathbf s_k\) represents the spatial pattern of a source group. The
% goal is to estimate a nonnegative source activity vector
% \[
% \boldsymbol\theta\in\mathbb R_+^K
% \]
% from sparse sensor observations. Known inventories make the problem
% lower-dimensional and scientifically meaningful, but they do not by
% themselves make the source groups identifiable.

% Wind introduces an additional complication: source contributions are not
% instantaneous. Emissions released at earlier times may become visible at
% a sensor only after being transported by the intervening wind field. We
% therefore model source-to-sensor influence through a lagged
% wind-conditioned response matrix. At sensor time \(t\), this response
% takes the form
% \[
% H_t^{\mathrm{lag}}
% =
% \sum_{\ell\in\mathcal L_t}
% \mathcal O
% G_{t,t-\ell}^{\partial}
% (\mathbf w_{t-\ell:t})
% S,
% \]
% where \(G_{t,t-\ell}^{\partial}\) maps emissions released at time
% \(t-\ell\) to concentrations at time \(t\) under open-boundary transport,
% and \(\mathcal O\) samples the resulting field at the sensor locations.
% Stacking over time gives a finite-horizon response matrix
% \[
% H_{1:T}^{\mathrm{lag}}
% \in
% \mathbb R^{mT\times K}.
% \]

% Real sensor observations also contain background and temporal components
% that are not explained by the local source inventory. We represent these
% components using a low-dimensional basis \(Q\), leading to the model
% \[
% \mathbf Y
% =
% H_{1:T}^{\mathrm{lag}}\boldsymbol\theta
% +
% Q\boldsymbol\gamma
% +
% \mathbf E.
% \]
% After projecting out the background subspace, the relevant inverse
% problem becomes
% \[
% P_Q^\perp \mathbf Y
% =
% P_Q^\perp H_{1:T}^{\mathrm{lag}}\boldsymbol\theta
% +
% P_Q^\perp \mathbf E.
% \]
% Thus, the central object in this work is the projected lagged
% source-response matrix
% \[
% \widetilde H
% =
% P_Q^\perp H_{1:T}^{\mathrm{lag}}.
% \]

% Our main claim is that the identifiable resolution of sparse-sensor
% source apportionment is determined by the geometry of
% \(\widetilde H\). If the columns of \(\widetilde H\) are linearly
% independent and well-conditioned, then source groups can be separated
% robustly. If the matrix is rank-deficient, ill-conditioned, weakly
% visible, or contains highly coherent source fingerprints, then some
% source groups cannot be reliably interpreted separately. In such cases,
% the data may support only a coarser attribution, such as the combined
% contribution of several ambiguous source groups.

% This perspective reframes source apportionment as an
% identifiability-aware reporting problem. The goal is not simply to return
% the finest possible vector of source contributions, but to determine
% which source contributions are defensible from the available sensors,
% wind fields, source inventories, background correction, and noise level.
% Accordingly, an apportionment method should report not only
% \(\widehat{\boldsymbol\theta}\), but also diagnostics such as rank,
% smallest singular value, condition number, effective rank, source
% visibility, pairwise coherence, and source groups that should be merged.

% We make the following contributions:
% \begin{itemize}[itemsep=0pt,topsep=0pt,leftmargin=*]
%     \item We formulate sparse-sensor source apportionment as a lagged
%     wind-conditioned inverse problem over known or proxy source
%     inventories.

%     \item We show that source identifiability is governed by the
%     projected lagged response matrix
%     \(\widetilde H=P_Q^\perp H_{1:T}^{\mathrm{lag}}\), rather than by
%     the raw source inventories or raw sensor fit alone.

%     \item We derive identifiability diagnostics based on rank,
%     \(\sigma_{\min}\), condition number, effective rank, source
%     visibility, background absorption, and pairwise source-fingerprint
%     coherence.

%     \item We propose an identifiability-aware source apportionment
%     framework that estimates source contributions while reporting
%     uncertainty, ambiguity, and source groups that should be merged.

%     \item We validate the framework in controlled city-scale pollution
%     simulations, testing whether attribution accuracy tracks the
%     predicted identifiability structure under varying noise levels, wind
%     regimes, background correction, and sensor geometries.
% \end{itemize}

% Overall, this work argues that source apportionment from sparse
% air-quality sensors is not merely a modeling or optimization problem. It
% is constrained by the information content of the sensing system. Reliable
% source attribution requires understanding which source fingerprints are
% actually separable after transport, lag, background correction, and
% noise.

\section{Introduction}
\label{sec:introduction}

Urban air pollution is driven by local emissions, meteorological transport,
regional background pollution, and sensor noise.  For policy, the central
question is often not only \emph{where} pollution is high, but \emph{which}
source groups contributed to the observed concentrations.  Source
apportionment therefore seeks to estimate the contributions of categories such
as traffic, industry, brick kilns, road dust, and population-related activity
from a combination of source inventories, transport models, and concentration
measurements.

This task is especially difficult in sparse sensor networks.  A model may fit
sensor trajectories accurately while still failing to identify the underlying
source contributions.  The reason is geometric: after wind transport and sparse
observation, two distinct source groups can induce nearly identical
sensor-time signatures.  A strong distant source and a weaker nearby source,
for example, may become observationally equivalent once transport attenuation,
finite travel time, background correction, and noise are taken into account.
In such cases, fine-grained source attribution is not merely statistically
uncertain; it is not identifiable from the available sensing system.

We study this issue in the inventory-based setting.  Instead of attempting to
recover an arbitrary unknown source field, we assume access to a finite set of
known or proxy source maps.  These inventories restrict attribution to a
lower-dimensional source model and make the problem interpretable.  To allow
realistic time variation without estimating one free value at every hour, we
represent each source activity as a nonnegative combination of a small number
of temporal basis functions.  The unknowns are therefore source--basis
coefficients, from which source activity trajectories and contributions are
reconstructed.  However, the inventory and temporal-basis restrictions alone
do not make the problem identifiable.  Coefficient fingerprints can still be
weakly visible, absorbed by background components, or mutually coherent after
transport to the sensor network.

The key object in our analysis is a projected lagged source-response matrix.
For a fixed wind sequence, transport operator, source inventory, sensor layout,
and lag window, each source--basis pair induces a finite-horizon sensor-space
fingerprint.  Stacking these fingerprints gives a lagged source--basis response
matrix \(H_{\Phi}^{\mathrm{lag}}\).  Real observations also contain background and
temporal components, represented by a low-dimensional basis \(Q\).  After
projecting out this background space, the inverse problem is governed by
\[
\widetilde H_\Phi = P_Q^\perp H_{\Phi}^{\mathrm{lag}},
\]
where \(P_Q^\perp\) is the projection onto the orthogonal complement of the
background basis.  The rank, singular values, coefficient visibility,
background absorption, pairwise coherence, and ray distance of
\(\widetilde H_\Phi\) determine which temporal source contributions can be
separated and which should be reported only after merging or qualification.
The familiar constant-activity model is recovered by using one constant basis
function per source.

This perspective changes the role of source apportionment models.  The goal is
not to return the most detailed vector of source contributions permitted by the
inventory.  It is to report the finest source resolution that is defensible
under the available sensors, wind regime, background correction, and noise
level.  An apportionment method should therefore output source estimates
alongside identifiability diagnostics and ambiguity certificates.

We make the following contributions:
\begin{itemize}[leftmargin=*,itemsep=0pt,topsep=0pt]
    \item We formulate sparse-sensor, inventory-based source apportionment as a
    lagged wind-conditioned inverse problem over nonnegative source--temporal-
    basis coefficients.
    \item We identify the projected lagged response matrix
    \(\widetilde H_\Phi=P_Q^\perp H_{\Phi}^{\mathrm{lag}}\) as the central
    object governing coefficient separability after transport, sparse
    observation, lag, and background correction.
    \item We derive diagnostics for exact and noise-robust identifiability,
    including rank, smallest singular value, condition number, effective rank,
    coefficient visibility, background absorption, pairwise fingerprint
    coherence, and ray distance.
    \item We propose an identifiability-aware apportionment procedure that fits
    nonnegative source--basis coefficients, reconstructs source activity
    trajectories, reports uncertainty, flags weakly visible coefficients, and
    merges source groups whose fingerprints are indistinguishable under the
    current sensing system.
    \item We design controlled city-scale advection--diffusion experiments to
    test whether source attribution error, source ambiguity, and the benefit of
    merging are predicted by the geometry of \(\widetilde H_\Phi\) under varying
    noise levels, wind regimes, background bases, and sensor layouts.
\end{itemize}

The broader message is that sparse-sensor source apportionment is limited by
the information content of the sensing system.  Reliable attribution requires
not only a transport model and an optimizer, but an explicit account of which
source--basis fingerprints are visible and separable in the observed sensor
space.
